\documentclass[10pt,a4paper]{article}
\usepackage[utf8]{inputenc}
\usepackage{amsmath}
\usepackage{amsfonts}
\usepackage{amssymb}
\usepackage{german}
\usepackage{graphicx}
\author{SEP - ITS - Team}
\title{Pflichtenheft}
\begin{document}
	\maketitle
	\includegraphics[scale=0.5, origin=c]{logo.png}
	\newpage
	
	\section{Zielbestimmung}
		Ziel des Projekts ist die autonome Erkundung eines unbekannten rechtwinkligen Spielfelds durch bis zu
		sechs e-puck Roboter. Dies wird durch Kooperation der Roboter und geeignete Algorithmen möglichst
		effizient umgesetzt. Zusätzlich erfolgt die kontinuierliche Visualisierung der Position der e-pucks auf der
		bereits erkundeten Karte mit einem Android-Smartphone. Die Steuerung eines Roboters kann wahlweise
		automatisch oder manuell per Smartphone erfolgen.
		\subsection{Musskriterien}
			\begin{itemize}
				\item e-puck Roboter
				\begin{itemize}
					\item Multiple Anzahl von Robotern
						\\ \textsl{Die Unterstützung von bis zu sechs e-puck Roboter muss gewährleistet sein.}
					\item Kommunikation über Bluetooth Netzwerk
					\item Erkundung einer unbekannten rechtwinkligen Fläche
						\\ \textsl{Das Spielfeld soll hierbei aus quadratischen Feldern bestehen. Diese sind rasterförmig
						 	und zusammenhängend angeordnet. Die Startpositionen der e-puck Roboter ist fest
						 	vordefiniert.}
					\item Fortbewegung auf den Kanten der Quadrate
						\\ \textsl{Die Linien müssen schwarz mit ausreichendem Kontrastverhältnis zum Untergrund
							sein.	Zwingend erforderlich ist, dass die Linienbreite innerhalb der Spezifikation der
							Bodensensoren	liegt.}		
					\item Kollisionsvermeidung
						\\ \textsl{Während der Erkundungsphase müssen die Roboter Kollisionen mit anderen Robotern
							bzw. Hindernissen vermeiden.}	
					\item Gleichberechtigung der Teilnehmer
						\\ \textsl{Sämtliche Roboter haben den selben Aufgabenbereich, auf eine Master/Slave-
							Beziehung wird verzichtet.}	
					\item Rückkehr zum Ausgangspunkt
						\\ \textsl{Nach vollendeter Erkundung müssen die e-pucks zurück zu ihrem jeweiligem
							Ausgangspunkt fahren.}	
				\end{itemize}
				\item Smartphone
				\begin{itemize}
					\item Kommunikation über Bluetooth
						\\ \textsl{Die Verbindung des Smartphones mit den Robotern muss über die vorhandenen
							Bluetooth Schnittstellen erfolgen.}
					\item Visualisierung
						\\ \textsl{Die bereits erkundete Karte muss inkl. der aktuellen Roboterpositionen in einer
							benutzerfreundlichen Android-Anwendung übersichtlich dargestellt werden.}		
					\item Steuerung des Roboters
						\\ \textsl{Es muss ein e-puck mit Hilfe der Anwendung gewählt werden können. Dieser lässt sich
							durch den Benutzer entlang der Kanten steuern. Die Richtung sowie Geschwindigkeit ist
							stufenweise anpassbar.}		
					\item Zwei Steuerungsarten
						\\ \textsl{Die Steuerung der Roboter erfolgt wahlweise über einen On-Screen- Joystick oder über
							den integrierten Beschleunigungssensor.}											
				\end{itemize}
			\end{itemize}
		\subsection{Wunschkriterien}
			\begin{itemize}
				\item e-puck Roboter
				\begin{itemize}
					\item Beliebige Startpositionen
						\\ \textsl{Die Roboter können auf frei wählbaren Startpositionen innerhalb des Spielfeldes
							abgesetzt werden.}
					\item Kritische Bereiche
					\item Virtuelle Tankstelle(n)
					\item Zustandsvisualisierung	
					\item Einstellbarer Synchronisationsmodus
				\end{itemize}
				\item Smartphone
				\begin{itemize}
					\item Exportfunktion
					\item Internationalisierung	
					\item Anzeige von Statusinformationen
					\item Pfadanzeige von einzelnen Robotern										
				\end{itemize}
			\end{itemize}
		\subsection{Abgrenzungskriterien}
			\begin{itemize}
				\item Keine abweichenden Spielfelder
					\\ \textsl{Die Felder und zugehörigen Linien müssen den genannten Musskriterien entsprechen.
						Es sind keine runden oder diagonalen Verbindungslinien erlaubt.}
				\item Keine dynamischen Änderungen
					\\ \textsl{Nachdem der Erkundungsvorgang gestartet wurde, dürfen keine Modifikationen am
						Spielfeld getätigt werden, welche Einfluss auf den Erkundungsalgorithmus zur Folge haben.}
				\item Keine dynamischen Änderungen
					\\ \textsl{Nachdem der Erkundungsvorgang gestartet wurde, dürfen keine Modifikationen am
						Spielfeld getätigt werden, welche Einfluss auf den Erkundungsalgorithmus zur Folge haben.}
				\item Maximal ein Smartphone
					\\ \textsl{Es darf lediglich ein Smartphone zur Auswahl, Steuerung und Visualisierung verwendet
						werden.}					
				\item Maximal ein Smartphone
					\\ \textsl{Es darf lediglich ein Smartphone zur Auswahl, Steuerung und Visualisierung verwendet
						werden.}		
			\end{itemize}
	\section{Produkteinsatz}
		\subsection{Anwendungsbereiche}
			Das Projekt dient der Grundlagenforschung in den Bereichen:
			\begin{itemize}
				\item Robotik
				\item Verteilte Systeme
				\item Künstliche Intelligenz
			\end{itemize}
		\subsection{Zielgruppen}
			\begin{itemize}
				\item Studenten
				\item Forschungsgruppen in ähnlichen Bereichen
				\item e-puck Community
			\end{itemize}
		\subsection{Betriebsbedingungen}
			\begin{itemize}
				\item Ausreichende Stromversorgung
					\\ \textsl{ Die Akkuleistung der einzelnen Roboter sowie des Smartphones muss für die gesamte
						Erkundungsdauer ausreichend sein.}
				\item Geeignete Bedingungen für Funknetzwerke
					\\ \textsl{ Die Ausmaße des Spielfelds, die Abstände von e-pucks oder des Smartphones sowie
						Signale anderer Netze dürfen keine störenden Einflüsse auf die Bluetooth-Verbindungen haben.}
			\end{itemize}
\end{document}